% ===================================================================
\chapter{The Baseline in Detail: Kiem's Thesis}
\label{ch:kiem}

Almost every number in this thesis is eventually compared with a number from
Lukas Kiem's master's thesis \cite{kiem2025}. A comparison is only as good as
one's understanding of what the other number measures, so this chapter sets out
what Kiem did, on what data, with what tools, and what he concluded --- in
enough detail that the later comparisons (Chapters~\ref{ch:replication}
and~\ref{ch:kiemcompare}) can be followed without his thesis open. Where
this project's early understanding of his work was wrong, the chapter says so,
because several of those misunderstandings shaped decisions described in
Chapter~\ref{ch:journey}.

\section{The problem he set himself}
\label{sec:kiem_problem}

Kiem worked at Infineon on the transport problem of
Section~\ref{sec:motivation}: a single high-resolution radar MMIC produces more
raw data than a cheap in-vehicle network can carry. He defined it
quantitatively with a use case (his Section~3.3,
Table~\ref{tab:kiem_usecase}). One Infineon CTRX8191F MMIC, four receive
channels, 12-bit IF samples at an output sample rate of \SI{50}{MS/s}, gives
\begin{equation}
50\,\mathrm{MS/s} \times 12\,\mathrm{bit} \times 4 = \SI{2.4}{Gbit/s},
\end{equation}
so carrying it over one 1000BASE-T1 Gigabit Ethernet link needs a compression
ratio of \num{2.4}, or \num{2.7} once protocol headers are counted. The
sensor cycles at \SI{20}{\hertz} (\SI{50}{\milli\second}); a frame of 1024
ramps of 1024 samples takes about \SI{21}{\milli\second} to acquire, and he
allowed compression \SI{15}{\milli\second} of the remaining
\SI{29}{\milli\second}.

\begin{table}[htbp]
\centering
\caption[Kiem's use case]{Kiem's use case (his Table~3.3) and the
requirements he derived from it (his Section~3.3).}
\label{tab:kiem_usecase}
\small
\begin{tabular}{@{}lr@{}}
\toprule
Parameter & Value\\
\midrule
MMICs / transmitters used & 1 / 1\\
Samples per ramp & 1024\\
Ramps per frame & 1024\\
Receive channels & 4\\
IF data width & 12 bits\\
Output sample rate & \SI{50}{MS/s}\\
Raw data rate & \SI{2.4}{Gbit/s}\\
Cycle time & \SI{50}{\milli\second} (\SI{20}{\hertz})\\
\midrule
Required compression ratio & 2.4 (2.7 with headers)\\
Compression time budget & \SI{15}{\milli\second}\\
\bottomrule
\end{tabular}
\end{table}

Two features of this setting matter for everything that follows, and both
differ from the sensor used in this thesis.

\paragraph{One transmitter.} Kiem's radar is a single MMIC operated with one
transmitter (his Table~4.4 fixes ``Number of TX: 1''). Consecutive ramps are
therefore consecutive observations by the same antenna pair, one pulse interval
apart. There is no transmitter interleaving on his ramp axis, so the question
this thesis turns on (Section~\ref{sec:tdm}) never arises for him.

\paragraph{Real-valued samples.} The CTRX delivers real IF samples. The FFT of
a real sequence is conjugate-symmetric, so Kiem keeps the first half of the
spectrum because the second half is redundant. The ColoRadar cascade delivers
complex $I/Q$ samples; its spectrum is not symmetric, and the halving applied
here is a deliberate choice to keep positive ranges
(Section~\ref{sec:refchain}), not a consequence of symmetry.

\section{His simulation workflow}
\label{sec:kiem_workflow}

Kiem evaluated compression schemes in a MATLAB environment (his Figure~3.1,
Section~3.2.1) whose stages this thesis adopted almost unchanged:

\begin{enumerate}
\item \textbf{Input data}: a cube of raw ADC samples, real-world or synthetic.
\item \textbf{Pre-processing}: shift the ADC samples left by
$16 - \mathit{bitWidth} + 16$ bits into a 32-bit full-scale word, then apply a
Hanning window of the ramp length.
\item \textbf{Range FFT}, divided by the number of samples (to prevent overflow
in his FFT IP core), keeping the first half.
\item \textbf{Compress} and \textbf{decompress}, including any quantisation a
scheme needs.
\item \textbf{Second-stage processing} in double precision: a Hanning-windowed
Doppler FFT divided by the number of ramps, non-coherent integration (NCI) of
the power over all receive channels, and peak detection combining local maxima
with a noise-based global threshold.
\item \textbf{Comparator}: detections and metrics of each scheme against the
reference path.
\end{enumerate}

His metrics (his Section~3.4.2) were the compression ratio; false negatives
and false positives against the reference; the noise floor estimate, defined
as ``the mean of the noisy part of the signal in dB'' in dBFS; and the SNR,
defined as the RMS signal of the detected peaks in dB minus the noise floor.
The reference format was 16-bit floating point (FP16). He does not state how
0\,dBFS is defined, what exactly the ``noisy part'' is, or the detector's
threshold. Section~\ref{sec:kr_assump} records how this thesis filled each gap.

\section{His data}
\label{sec:kiem_data}

Kiem used two kinds of input (his Section~3.2.2), and it is important to be
precise about which produced which result, because this project's early
documents were not.

\paragraph{Real-world recordings.} Measurements taken with the CTRX in several
scenarios, from a single stationary target to several moving ones, all with the
configuration of Table~\ref{tab:kiem_data}: one MMIC, 512 samples, 512 ramps,
four receive channels, 14-bit IF data. His design-space table (Table~3.5), his
removal of run-length coding (Table~3.6), the derivation of his fixed Huffman
dictionary (``the real radar data with id~=~1'', Section~3.5.2) and his hardware
evaluation (the SD-card input of Section~5.2.1) all use these recordings. They
are not publicly available.

\paragraph{Synthetic scenes.} Generated from a list of targets (range bin,
Doppler bin, amplitude) with additive white Gaussian noise, for configurations
where no recording existed. He shows one such scene (Table~\ref{tab:kiem_data},
his Table~3.2 and Figures~3.5--3.7) but reports no compression results on it.

\begin{table}[htbp]
\centering
\caption[Kiem's data]{Kiem's two kinds of input data (his Tables~3.1
and~3.2).}
\label{tab:kiem_data}
\small
\begin{tabular}{@{}lll@{}}
\toprule
& Real-world recordings (Table 3.1) & Synthetic scene (Table 3.2)\\
\midrule
MMICs & 1 & 1\\
Samples $\times$ ramps & $512 \times 512$ & $1024 \times 512$\\
Receive channels & 4 & 4\\
IF bit width & 14 & 14\\
Targets & real scenes & 5: range bins [100, 150, 150, 200, 400],\\
& & Doppler bins [50, 50, 100, 200, 250],\\
& & amplitudes [0.8, 0.3, 0.4, 0.5, 0.7]\\
Noise & real & AWGN, $-75$\,dBFS after NCI\\
Used for & every published result & illustration only\\
\bottomrule
\end{tabular}
\end{table}

\begin{keybox}
\textbf{A correction to this project's own record.} The Thesis~A report
stated that Kiem's FPGA implementation ``was tested exclusively using simulated
dummy data'', and the executive summary that he ``evaluated compression
algorithms using simulated data with AWGN''. Both are wrong. Every compression
result Kiem publishes, in simulation and in hardware, was measured on his
real CTRX recordings; his synthetic scenes appear only as illustrations. What
is true is that his recordings are private, come from a single-transmitter
sensor, and are much smaller than a cascade frame --- which is the real case for
validating the algorithm again on public MIMO data, and is how the first
objective of this thesis is now stated (Section~\ref{sec:objectives}).
\end{keybox}

\section{His design space}
\label{sec:kiem_space}

Kiem surveyed the automotive schemes available to him (his Section~2.7) and
took all but one into a design-space exploration (his Table~3.4).

\begin{description}[leftmargin=!,labelwidth=24mm,style=nextline,font=\normalfont\bfseries]
\item[FP16] Half-precision floating point: the reference format.
\item[CMPRA, CMPRB] Infineon AURIX common-exponent formats that pack 16 or 32
complex bins into one 256-bit word with a shared exponent: a fixed CR of 2 or 4
against 16-bit data.
\item[EGE10 \ldots EGE4] Texas Instruments' compression engine: order-$k$
exponential-Golomb coding of blocks of range bins, with least significant bits
dropped to meet a fixed CR of 16/$k$.
\item[DRHE] Doppler Redundancy Hybrid Encoding, from Liang Li's master's thesis
\cite{li2014} and Zivkovic and Li \cite{zivkovic2015}: prediction along the
Doppler (slow-time) dimension of the range-FFT output with two IIR filters in
polar form, then run-length and Huffman coding of the differences. The
original papers report ``ideal'' compression ratios of 9.7--9.8 (his
Table~2.1).
\end{description}

\noindent
The exception was \textbf{LPC-Huffman}, the method of his reference [25]
(Meucci and Mancuso \cite{meucci2022}). Its treatment is quoted here in full,
because Section~\ref{sec:lpc_notkiem} and Chapter~\ref{ch:replication} depend
on exactly what it says.

\begin{quote}\small
``In the work [25], the authors present a solution to compress FMCW radar data
using a combination of linear predictive coding (LPC) and Huffman coding. The
input data is said to be int16 data. Therefore, we assume that this work is
performed on raw ADC data.'' (Section~2.7.4)

``The test results in the work achieve a CR of 2.056.'' (Section~2.7.4)

``We take all these solutions except the LPC-Huffman compression algorithm into
the design space. LPC-Huffman compression is not taken because it has the same
approach as DRHE, but instead of compressing range FFT data, it operates on raw
ADC data; the results for DRHE are more promising.'' (Section~3.4.1)
\end{quote}

\noindent
Three things follow directly from these sentences. Kiem \emph{assumed} the
domain of \cite{meucci2022} rather than knowing it. He did not implement or run LPC-Huffman.
And ``more promising'' compares two published numbers --- 2.056 on the data of
\cite{meucci2022} against 9.7--9.8 on the data of the DRHE papers --- not two algorithms on
the same data. His own Table~3.4 describes DRHE as ``LPC on range FFT data'':
in his words the two methods differ in \emph{where} they predict, not in
\emph{whether}. Chapter~\ref{ch:replication} runs the method he excluded on
his own scene.

\section{His design-space result}
\label{sec:kiem_dse}

Table~\ref{tab:kiem_t35} reproduces his Table~3.5. On his real recordings DRHE
reached CR \num{3.25} with exactly FX16's detection metrics --- it is lossless
with respect to the 16-bit codes, so the only degradation against FP16 comes
from quantising to 16-bit fixed point. The lossy schemes either stayed at or
below the CR requirement or degraded detection markedly.

\begin{table}[htbp]
\centering
\caption[Kiem's design-space results]{Kiem's design-space exploration on his
real recordings (his Table~3.5; Table~5.1 repeats it). FP\% is normalised by
test detections, so it is a false discovery rate
(Section~\ref{sec:detectionmetrics}).}
\label{tab:kiem_t35}
\small
\begin{tabular}{@{}lrrrrr@{}}
\toprule
Algorithm & CR & FN\% & FP\% & NF (dBFS) & SNR (dB)\\
\midrule
FP16  & 1    & 0     & 0     & $-70.714$ & 22.407\\
FX16  & 1    & 0.59  & 1.18  & $-70.691$ & 22.330\\
CMPRA & 2    & 0.59  & 0.59  & $-70.687$ & 22.418\\
CMPRB & 4    & 7.69  & 4.73  & $-70.253$ & 22.111\\
EGE10 & 1.6  & 1.18  & 1.18  & $-70.689$ & 22.403\\
EGE8  & 2    & 2.37  & 2.37  & $-70.640$ & 22.289\\
EGE6  & 2.67 & 7.69  & 7.10  & $-70.367$ & 22.197\\
EGE4  & 4    & 40.24 & 26.63 & $-74.689$ & 28.245\\
\textbf{DRHE} & \textbf{3.25} & 0.59 & 1.18 & $-70.691$ & 22.330\\
\bottomrule
\end{tabular}
\end{table}

\noindent
A false-negative rate of \SI{0.59}{\percent} is one missed detection in about
170, which is a useful check on the scale of his scene: a single frame of his
recording contained on the order of 170 reference detections.

\section{How he adapted DRHE for hardware}
\label{sec:kiem_adapt}

Two simplifications turned the published DRHE into his final scheme.

\paragraph{Remove the run-length code (his Section~3.5.1).} The original coder
emits an 8-bit symbol RRRRSSSS per non-zero residual --- a 4-bit run of
preceding zeros and a 4-bit size category --- Huffman-codes the symbol and
appends the value. On his data the R4S4 histogram was far from the one the
dictionary had been designed for, and dropping the run-length half cost almost
nothing: CR \num{3.25} with R4S4 against \num{3.24} with S4 alone
(Table~\ref{tab:kiem_t36}), while shrinking the dictionary from 256 symbols to
16 and bounding the longest codeword.

\begin{table}[htbp]
\centering
\caption[Kiem: DRHE without RLE]{Kiem's Table~3.6: removing the run-length
part of the coder.}
\label{tab:kiem_t36}
\small
\begin{tabular}{@{}lrrrrr@{}}
\toprule
Algorithm & CR & FN\% & FP\% & NF (dBFS) & SNR (dB)\\
\midrule
DRHE (R4S4 + RLE) & 3.25 & 0.59 & 1.18 & $-70.691$ & 22.33\\
DRHE (S4 only)    & 3.24 & 0.59 & 1.18 & $-70.691$ & 22.33\\
\bottomrule
\end{tabular}
\end{table}

\paragraph{Fix the Huffman dictionary (his Section~3.5.2).} A dictionary built
per frame would need a histogram pass, a tree construction and transmission of
the table, none of which suit a streaming IP core. He therefore built one
dictionary once, from ``the real part of the S4 values for the real radar data
with id = 1'', and fixed it for simulation and hardware (his Appendix~A; this
thesis's Appendix~\ref{app:huffman}). Its shortest codewords --- two bits ---
go to categories 2, 3 and 4, i.e.\ residuals of magnitude 2 to 15, so the
dictionary is a record of what his residuals looked like.
Section~\ref{sec:kr_fingerprint} uses exactly that property.

\section{His final compression scheme}
\label{sec:kiem_final}

His Section~3.5.3 specifies the scheme implemented in hardware; this thesis
audits its own implementation against it line by line in
Section~\ref{sec:kc_audit}. In summary, for each range bin $n$ of ramp $m$ and
each receive channel:
\begin{align}
\hat{a}[m,n] &= \alpha\hat{a}[m-1,n] + (1-\alpha)\,a[m-1,n], \\
\hat{\theta}[m,n] &= \beta\hat{\theta}[m-1,n] + (2-\beta)\,\theta[m-1,n] - \theta[m-2,n],
\end{align}
with $\alpha = 0.6$, $\beta = 0.4$, $\hat\theta$ wrapped once into
$[-\pi,\pi]$, the prediction converted back to Cartesian with rounding, and the
difference $X - \hat{X}$ taken in 16-bit two's complement. Each difference is
coded as its size category S4 (Table~3.7 of his thesis, reproduced here as
Table~\ref{tab:catranges}), Huffman-coded with the fixed dictionary, followed by
S4 APPEND bits holding $v$ for positive and $v + 2^{S_4} - 1$ for negative
values. The model state is initialised to zero (his Section~3.5.4) and updated
from the input samples.

\section{His hardware}
\label{sec:kiem_hw}

He implemented the scheme with Vitis HLS and Vivado 2022.2 on an AMD Zynq
UltraScale+ \texttt{XCZU3CG-SFVC784-1-e} on a Trenz TE0820 module, inside a
complete system: a bare-metal program loads pre-windowed 16-bit ADC data from
an SD card, an AXI DMA streams it through a Xilinx FFT IP core into the
compression IP, and the compressed stream is written back to memory
(his Figure~4.2). The compression IP (his Section~4.3.4):

\begin{itemize}
\item reads one 128-bit AXI4-Stream word per cycle --- one complex 16-bit bin
for each of four channels --- and processes the channels in parallel;
\item uses fixed point throughout: magnitudes \texttt{ap\_ufixed<16,16>},
phases \texttt{ap\_fixed<16,3>}, intermediates \texttt{ap\_fixed<32,4>}, with
\texttt{hls::sqrt}, \texttt{hls::atan2}, \texttt{hls::sin} and
\texttt{hls::cos};
\item holds the five state arrays per channel in block RAM, $4 \times 1024$
words each;
\item packs each channel's two codes into a 58-bit word and the four words into
a 256-bit output buffer;
\item runs at $\mathrm{II} = 1$ at \SI{100}{\mega\hertz}.
\end{itemize}

\noindent
Table~\ref{tab:kiem_hwres} collects the numbers he reports.

\begin{table}[htbp]
\centering
\caption[Kiem's hardware results]{Kiem's published hardware figures
(Section~5.2 and Tables~5.2, 5.3, 5.5 of \cite{kiem2025}).}
\label{tab:kiem_hwres}
\small
\begin{tabular}{@{}lr@{}}
\toprule
Quantity & Value\\
\midrule
Compression IP: LUT / FF / BRAM (36\,Kb) / DSP & \num{45892} / \num{9243} / 10 / 44\\
Whole system: LUT / FF / BRAM / DSP & \num{59353} / \num{32552} / 38 / 80\\
Clock & \SI{100}{\mega\hertz}\\
IP latency (HLS estimate) & 23 cycles, \SI{230}{\nano\second}\\
IP latency (logic analyser, to first output word) & 31 cycles\\
System latency (last DMA read to last DMA write) & \SI{79410}{\nano\second}\\
Input throughput & ``\SI{11.92}{Gbit/s}''\\
CR, MATLAB (double-precision FFT) & 3.25\\
CR, bit-accurate simulation & 3.24\\
CR, hardware & 3.17\\
FN / FP after hardware FFT & \SI{1.18}{\percent} / \SI{2.37}{\percent}\\
\bottomrule
\end{tabular}
\end{table}

\noindent
Chapter~\ref{ch:interpret} shows that the throughput is a theoretical figure
in binary units ($128 \times \SI{100}{\mega\hertz} = \SI{12.8}{Gbit/s} =
\SI{11.92}{\gibi bit/s}$), and that the block RAM is counted in 36\,Kb tiles.
His hardware CR is lower than his simulation CR because the hardware FFT is
fixed point: the same ADC data produces slightly different 16-bit codes, which
he confirmed by simulating the FFT core's bit-accurate model.

\section{What he concluded, and what he left open}
\label{sec:kiem_open}

He concluded that DRHE met every requirement: CR \num{3.17} against 2.4 (2.7
with headers), \SI{0.08}{\milli\second} against a \SI{15}{\milli\second}
budget, and ``99.4\,\% detection matches with the reference''. His
limitations and future work (his Chapter~6) were:

\begin{itemize}
\item the precision lost by 16-bit fixed point before and in the compressor,
with the suggestion that, since the IP core far exceeds its real-time
requirement, \emph{floating-point} arithmetic for the FFT and the model
prediction could be investigated;
\item a bit-accurate decompression model;
\item tests with a real radar sensor in the loop and all configuration options;
\item safe error handling and a guaranteed compression ratio;
\item ``transpose compression'' to suit the Doppler FFT's memory access.
\end{itemize}

\noindent
This thesis took up the first two directly: its DRHE uses floating-point
prediction (Section~\ref{sec:drheimpl}; the cost is measured in
Chapter~\ref{ch:interpret}), and it implements and verifies a decompressor for
every design. It could not take up the third, since no board was available,
and did not attempt the last two.

\section{What the later chapters take from this one}

\begin{enumerate}
\item Kiem's published compression ratios come from private recordings of a
single-transmitter, real-sampled radar. No number in this thesis was measured
on the same data, so a CR here and a CR there are not like-for-like until
something connects them. Chapter~\ref{ch:replication} is that connection.
\item Kiem did not measure LPC. His statement that DRHE is ``more promising''
compares published figures from different data, for a method (fast-time
prediction of raw samples) that is not the LPC of this thesis.
\item His final DRHE scheme is specified in enough detail to implement exactly,
and Section~\ref{sec:kc_audit} checks that this thesis did.
\item His fixed dictionary is a fingerprint of the residual statistics of his
own data, which Section~\ref{sec:kr_fingerprint} reads.
\end{enumerate}
